\section{Client Plugin}
\label{sec:ClientPlugin}

The client plugin is an Eclipse RCP bundle that runs inside Archi. It captures
local EMF model changes, submits them to the Archimesh server, and applies
remote operations into the local model.

\subsection{Plugin entry point}

The plugin activator is \\ \texttt{org.gautelis.archimesh.plugin.ArchimeshPlugin},
an \texttt{AbstractUIPlugin} subclass.

On startup it creates the three central objects:

\begin{description}
\item[\texttt{ArchimeshSessionManager}] WebSocket lifecycle, outbox handling,
  join/rejoin decisions and local submit tracking.
\item[\texttt{WorkbenchModelLifecycleBridge}] Hooks into Archi's workbench to
  detect model open, close, save and dirty-state changes.
\item[\texttt{ArchimeshStatusLineBridge}] Displays connection state, conflict
  banners and outbox capacity warnings in the Archi status line.
\end{description}

The plugin id is \texttt{org.gautelis.archimesh.plugin}. Debug and trace
logging are controlled by the system property \texttt{archimesh.debug} and
\texttt{archimesh.trace}, or equivalently the environment variables
\texttt{ARCHIMESH\_DEBUG} and \texttt{ARCHIMESH\_TRACE}.

\subsection{Session manager}

\texttt{ArchimeshSessionManager} is the central client-side coordinator. It
manages:

\begin{itemize}
\item WebSocket connection lifecycle (connect, disconnect, reconnect)
\item Per-model persistent bounded outbox (FIFO queue with replay)
\item Join/rejoin decision logic with cache metadata
\item Echo suppression for locally submitted op batches
\item Cache autosave policy (debounced saves on sync traffic)
\end{itemize}

\subsubsection{Connection flow}

\begin{enumerate}
\item The user selects \textbf{Tools \textgreater\ Connect Archimesh...} and
  provides the server URL, model id and credentials.
\item The session manager resolves the join decision:
  \begin{itemize}
  \item If the session is server-backed, cache metadata is read from
    \filepath{\textasciitilde/Archi/archimesh-cache/<modelId>.meta.properties}.
  \item If metadata is valid and the cached revision is known, the join
    includes \texttt{lastSeenRevision} so the server can send a delta.
  \item If metadata is missing, mismatched or the revision is unknown, the
    join omits \texttt{lastSeenRevision} to request a cold snapshot.
  \end{itemize}
\item The session manager opens a WebSocket to
  \texttt{<baseUrl>/models/<modelId>/stream} and sends a \texttt{Join} message.
\item The server responds with \texttt{CheckoutSnapshot} (full model state)
  or \texttt{CheckoutDelta} (incremental changes since lastSeenRevision).
\item The client applies the checkout payload into the local Archi model.
\end{enumerate}

\subsubsection{Rejoin and cache verification}

On reconnect the session manager reads the cache metadata file. A cache
projection is discarded (forcing a cold snapshot) if:

\begin{itemize}
\item The metadata file does not exist or cannot be read.
\item The metadata \texttt{serverBacked} flag does not match.
\item The \texttt{modelId} or \texttt{modelRef} does not match the target.
\item The cache file is missing from disk.
\item The cached revision is unknown (negative).
\item The cache file exists but is implausibly small for a model with
  revision \(\geq 1\) (under 1~KB --- likely truncated or corrupted).
\end{itemize}

After join, the server's first revision hint is compared against the cache
revision used at join. If the server revision is behind the cache revision,
the client requests a cold snapshot rebuild.

\subsection{EMF change capture}

\texttt{EmfChangeCapture} observes local EMF model changes via Archi's
ECore notification mechanism. When the user edits the model, the capture
layer:

\begin{enumerate}
\item Receives EMF notifications for element, relationship, view, folder
  and property changes.
\item Batches correlated changes (e.g. a view object create + its initial
  geometry update) into a single logical operation.
\item Defers connection-creation submission until all referenced identifiers
  (source, target, relationship, view) are available, then emits the
  relationship and connection in a single \texttt{SubmitOps} batch.
\item Passes the collected operations to \texttt{OpMapper} for wire-format
  encoding.
\end{enumerate}

The capture layer runs under a remote-apply guard so that operations applied
from the server are not re-captured and re-submitted as local edits.

\subsection{Operation mapping}

\texttt{OpMapper} converts captured EMF changes into Archimesh wire-format
operations. The supported operation types are:

\begin{longtable}{p{52mm}p{102mm}}
\toprule
Operation & Purpose \\
\midrule
\texttt{CreateElement} & Create a new ArchiMate element. \\
\texttt{UpdateElement} & Update element name or documentation. \\
\texttt{DeleteElement} & Delete an element and cascade-delete its
  diagram representations. \\
\texttt{CreateRelationship} & Create a relationship between two elements. \\
\texttt{UpdateRelationship} & Update relationship fields. \\
\texttt{DeleteRelationship} & Delete a relationship and cascade-delete
  its connections. \\
\texttt{CreateFolder} & Create a user folder. \\
\texttt{UpdateFolder} & Rename a folder or update its documentation. \\
\texttt{DeleteFolder} & Delete an empty user folder. \\
\texttt{CreateView} & Create a new diagram view. \\
\texttt{UpdateView} & Update view metadata. \\
\texttt{DeleteView} & Delete a view and its contents. \\
\texttt{CreateViewObject} & Place an element on a diagram. \\
\texttt{UpdateViewObjectOpaque} & Update geometry, style or label fields
  of a diagram object. \\
\texttt{DeleteViewObject} & Remove a diagram object. \\
\texttt{CreateConnection} & Draw a connection between two diagram objects. \\
\texttt{UpdateConnectionOpaque} & Update connection style, label or
  bendpoint fields. \\
\texttt{DeleteConnection} & Remove a diagram connection. \\
\texttt{AddElementToFolder} & Place an element in a folder. \\
\texttt{RemoveElementFromFolder} & Remove an element from a folder. \\
\texttt{SetProperty} & Set a key-value property on a model entity. \\
\texttt{UnsetProperty} & Remove a property. \\
\texttt{AddPropertySetMember} & Add an element to a property set. \\
\texttt{RemovePropertySetMember} & Remove an element from a property set. \\
\texttt{AddViewObjectChildMember} & Nest a view object inside another. \\
\texttt{RemoveViewObjectChildMember} & Remove nested view object membership. \\
\bottomrule
\end{longtable}

Each operation carries a \texttt{causal} object with \texttt{clientId},
\texttt{lamport} and \texttt{opId} for LWW merge semantics.

\subsection{Remote operation application}

\texttt{RemoteOpApplier} applies snapshots and remote broadcasts into the
local EMF model. Key behaviors:

\begin{itemize}
\item Applies operations under a remote-apply guard to prevent re-capture.
\item Applies per-field LWW merges for notation fields (geometry, style,
  labels) using causal clock comparison.
\item Maintains element field clocks and tombstones for delete safety.
\item Handles property-level LWW and OR-set semantics for
  \texttt{SetProperty}, \texttt{UnsetProperty},
  \texttt{AddPropertySetMember}, \texttt{RemovePropertySetMember}.
\item Restores folder structure before folder memberships during snapshot
  application.
\item Restores views before view objects and connections.
\item Defers dependency-sensitive operations when referenced entities are not
  yet created.
\end{itemize}

\subsection{Offline outbox}

When the WebSocket is unavailable or a live send fails, the client queues
operations in a persistent bounded outbox. Key guarantees:

\begin{itemize}
\item \textbf{FIFO per model.} Operations are drained in submission order.
\item \textbf{Rebase before send.} Each replayed op has its
  \texttt{baseRevision} updated to the current known server revision.
\item \textbf{Ack-gated removal.} Entries are removed only after the server
  sends a matching \texttt{OpsAccepted}. Websocket send completion is not
  enough --- the server must confirm commit.
\item \textbf{Deterministic retry.} Failed sends are retried with capped
  exponential backoff. Entries failing after the maximum replay attempts (5)
  are dropped as poison entries.
\item \textbf{Conflict resolution.} Entries receiving
  \texttt{PRECONDITION\_FAILED} or \texttt{LOCK\_CONFLICT} from the server
  are deterministically dropped. A conflict callback is fired, surfaced in
  the UI status line.
\item \textbf{Bounded queue.} Maximum 1000 entries. When full, the oldest
  entry is dropped. A warning fires at 80\% capacity (800 entries) and is
  displayed in the status line.
\item \textbf{Durable persistence.} The outbox is persisted to
  \filepath{\textasciitilde/Archi/archimesh-cache/<modelId>.archimate.outbox.properties}
  and reloaded on reconnect.
\end{itemize}

\subsection{Cache policy}

When running in server-backed mode, the client treats the local
\texttt{.archimate} file as a cache projection, not the primary artifact.

\begin{description}
\item[Cache path] \filepath{\textasciitilde/Archi/archimesh-cache/<modelId>.archimate}
\item[Metadata file] \filepath{<cacheFile>.meta.properties} stores
  \texttt{modelId}, \texttt{modelRef}, \texttt{lastKnownRevision},
  \texttt{serverBacked} and \texttt{savedAtEpochMs}.
\item[Autosave triggers] Debounced saves (2 s cooldown) after
  \texttt{CheckoutSnapshot}, \texttt{CheckoutDelta},
  \texttt{OpsBroadcast} and \texttt{OpsAccepted}.
\item[Force-save triggers] On model attach, detach and disconnect.
\item[Discard policy] Cache is discarded if inconsistent, invalid or not
  server-backed at reconnect.
\end{description}

\subsection{Echo suppression}

To avoid double-applying their own submitted operations, clients track the
most recent 256 locally submitted \texttt{opBatchId} values. When an
\texttt{OpsBroadcast} arrives via the Kafka fan-out path, the client checks
whether the batch was locally submitted. If so, the broadcast is silently
ignored. The local model was already updated optimistically and will be
corrected by the \texttt{OpsAccepted} confirmation if needed.

\subsection{Menu commands}

The plugin contributes the following commands to Archi's \textbf{Tools} menu:

\begin{description}
\item[Connect Archimesh...] Connect the active model to an Archimesh server.
  Opens a dialog for server URL, model id, user, session id and optional
  bearer token.
\item[Disconnect Archimesh] Disconnect from the server, save cache and
  return to local-only mode.
\item[Open Archimesh Model from Server...] Browse the server model catalog
  and open a model from the server.
\item[Switch Active Model to Server-Backed...] Attach Archimesh to the
  currently active local model, making it server-backed.
\item[Show Last Archimesh Conflict] Display details for the most recent
  conflict (precondition failure or lock conflict).
\item[Resynchronize Model] Discard the local projection and pull a fresh
  snapshot from the server.
\end{description}
